\documentclass[11pt,a4paper,twoside]{book}
% Shared parent preamble for the PhD thesis.
% Loaded by main.tex, and by every chapters/*.tex subfile via
% \documentclass[../main.tex]{subfiles}.

% --- Core math ---
\usepackage{amsmath,amssymb,amsfonts,mathtools,mathrsfs,dsfont,cancel,bm}
\usepackage{upgreek}
\usepackage{nicefrac}
\usepackage{commath}
\usepackage{braket}
\usepackage[nointegrals]{wasysym}
\usepackage{pifont}
\usepackage{ar}           % supplies \AR, the aspect ratio N/f in internal_tide_eddy
% Italic capital Xi, used in internal_tide_eddy for the placeholder field whose
% flux fraction is computed. No package loaded here defines it, so it is defined
% below. \providecommand steps aside if some future package supplies its own.
\providecommand{\varXi}{\mathit{\Xi}}

% --- Figures, tables ---
\usepackage{graphicx}

\usepackage[export]{adjustbox}
\usepackage{subcaption}
\usepackage{wrapfig}
\usepackage{float}
\usepackage{booktabs}
\usepackage{longtable}  % the list of symbols runs over several pages
\usepackage{dcolumn}
\usepackage{tikz,pgf,import}
\DeclareCaptionJustification{justified}{\justifying}
\captionsetup{justification=justified,singlelinecheck=false,labelfont=large}
\captionsetup[subfigure]{position=top,belowskip=0pt,justification=centering}

% --- Colors, highlighting, track changes ---
\usepackage[dvipsnames]{xcolor}
\usepackage{soul}
\usepackage[normalem]{ulem}
% \usepackage[inline]{trackchanges}

% --- Layout ---
% Wide text block, as in the reference thesis: the default book-class
% margins leave the text only ~121mm wide on A4, this widens it to ~160mm.
% headheight holds two lines, so a long chapter title in the running header
% cannot run into the rule above the text.
\usepackage[a4paper,inner=25mm,outer=25mm,top=25mm,bottom=30mm,headheight=26pt,headsep=12pt]{geometry}
\usepackage{ragged2e}
% No hyphenation anywhere. 10000 is TeX's infinity, so no word is ever split.
\hyphenpenalty=10000
\exhyphenpenalty=10000
\sloppy
\usepackage{parskip}
\usepackage{setspace}
\usepackage{microtype}
\usepackage{enumitem}

\setstretch{1.25}
\setlength{\parskip}{1.25em}
% --- Hyperlinks ---
% Loaded here, before orcidlink, because orcidlink silently \RequirePackages
% hyperref with no options if it isn't already loaded -- loading it later
% with options triggers "Option clash for package hyperref".
\usepackage[
    colorlinks=true,
    linkcolor=blue!60!black,
    citecolor=blue!60!black,
    urlcolor=blue!60!black
]{hyperref}
% \usepackage[nameinlink]{cleveref}

\usepackage{orcidlink}
\usepackage{subfiles}
\usepackage{xspace}  % used by the internal_tide_eddy chapter's \lr/\ir/\hr-family macros

% --- Bibliography ---
% One combined reference list at the end of the thesis. Resources are
% searched in the order listed, so the shared Zotero file is tried first
% and each paper's own .bib is the fallback for keys that predate it.
% APS-style citations: bracketed numbers, numbered in order of first use,
% consecutive runs compressed to [1-3].
\usepackage[backend=biber,style=numeric-comp,sorting=none,natbib=true]{biblatex}
\addbibresource{/Users/rajarshi/Dropbox/LaTeX/global_references_zotero.bib}
% Copies of each paper's own .bib, kept here because the paper folders
% themselves are not uploaded.
\addbibresource{chapters/bib_active_caustics.bib}
\addbibresource{chapters/bib_internal_tide_eddy.bib}
\addbibresource{chapters/bib_thermalization.bib}
\addbibresource{chapters/bib_inertial_caustics.bib}
\addbibresource{chapters/bib_polydisperse.bib}

% The entry title carries the link: to the DOI when the entry has one,
% otherwise to the URL. The DOI, URL and arXiv eprint are never printed
% as standalone fields.
\ExecuteBibliographyOptions{doi=false,url=false,eprint=false,isbn=false}
\newbibmacro*{linkedtitle}[1]{%
  \iffieldundef{doi}
    {\iffieldundef{url}
       {#1}
       {\href{\thefield{url}}{#1}}}
    {\href{https://doi.org/\thefield{doi}}{#1}}}
\DeclareFieldFormat{title}{\usebibmacro{linkedtitle}{\mkbibemph{#1}}}
\DeclareFieldFormat[article,inbook,incollection,inproceedings,patent,thesis,unpublished]%
  {title}{\usebibmacro{linkedtitle}{\mkbibquote{#1\isdot}}}

% --- Thesis accent color + chapter/section style ---
% Matches the grey "Chapter N" numeral with a bold title set below it,
% used as the visual reference for this thesis.
\definecolor{thesisaccent}{HTML}{8C3A2E}
\usepackage{titlesec}
% Chapter heads are set in TeX Gyre Adventor (family qag), a geometric sans
% whose round letters are much wider than the default one, at sizes a few
% steps above \Huge. Everything else keeps the ordinary sans.
\titleformat{\chapter}[display]
  {\normalfont\fontfamily{qag}\fontsize{36}{48}\selectfont\color{gray}}
  {\chaptertitlename\ \thechapter}
  {10pt}
  {\raggedleft\normalfont\sffamily\bfseries\Huge\color{black}}
\titleformat{\section}
  {\normalfont\sffamily\bfseries\fontsize{18}{20}\selectfont}
  {\thesection}{1em}{}
\titleformat{\subsection}
  {\normalfont\sffamily\bfseries\fontsize{14}{16}\selectfont}
  {\thesubsection}{0.75em}{}
\titleformat{\subsubsection}
  {\normalfont\sffamily\bfseries\itshape}
  {\thesubsubsection}{0.5em}{}

% --- Chapter summary box ---
% Opens each chapter in place of the paper's abstract. breakable lets a long
% summary run across a page break.
\usepackage[most]{tcolorbox}
\definecolor{clbox}{RGB}{100, 87, 166}
\tcbset{colframe = clbox}
\newtcolorbox{summary}{%
  breakable, enhanced, sharp corners,
  boxrule=0pt, colframe=clbox!6,
  borderline west={3pt}{0pt}{clbox},
  colback=clbox!6, colbacktitle=clbox!6, coltitle=clbox,
  title=Summary, fonttitle=\rmfamily\bfseries,
  fontupper=\small\itshape
}

% --- Running headers ---
\usepackage{fancyhdr}
\pagestyle{fancy}
\fancyhf{}
\fancyhead[LE]{\thepage\quad\nouppercase{\leftmark}}
\fancyhead[RO]{\nouppercase{\rightmark}\quad\thepage}
\renewcommand{\headrulewidth}{0.4pt}

\setlength{\parskip}{0.8em}
\setlength{\parindent}{0pt}

% ====================== Shared macros ======================
% Merged from Inertial_caustics/r_preamble.tex, polydisperse/r_preamble.tex,
% and the inline preambles of the Active Caustics, JPO and Thermalization
% manuscripts. Where two papers defined the same name differently, the
% more common / more recent version was kept.

\newcommand{\Dp}[3]{\frac{\partial^{#3} #1}{\partial #2^{#3}}}
\newcommand{\Dpc}[3]{\frac{\cancel{\partial}^{#3}#1}{\cancel{\partial} #2^{#3}}}
\newcommand{\dk}[2]{\frac{d^{#2} #1}{(2\pi)^{#2}}}
\renewcommand{\bra}[1]{\left<#1\right|}
\renewcommand{\ket}[1]{\left|#1\right>}
\renewcommand{\braket}[2]{\left<#1| #2\right>}
\newcommand{\mc}[1]{\mathcal{#1}}
\newcommand{\ms}[1]{\mathscr{#1}}
\renewcommand{\md}[1]{\mathds{#1}}
\newcommand{\mb}[1]{\mathbb{#1}}
\newcommand{\tx}[1]{\mbox{#1}}
\newcommand{\mtx}[1]{\scriptstyle #1 \textstyle}
\newcommand{\ttx}[1]{\scriptscriptstyle #1 \textstyle}
\newcommand{\eqn}{equation}
\newcommand{\soln}{solution}
\newcommand{\curl}[1]{\vec{\nabla} \times \vec{#1}}
\newcommand{\grad}[1]{\vec{\nabla}#1}
\newcommand{\lap}[1]{\nabla^2 #1}
\newcommand{\mean}[1]{\langle #1\rangle}
\newcommand{\eref}[1]{Eq.~\eqref{#1}}
\newcommand{\fref}[1]{Fig.~\ref{#1}}
\newcommand{\vect}[1]{\boldsymbol{#1}}

\renewcommand{\d}[3]{\frac{d^{#3} #1}{d #2^{#3}}}
\renewcommand{\t}[1]{\tilde{#1}}
\renewcommand{\vec}[1]{\boldsymbol{#1}}
\renewcommand{\implies}{\Rightarrow}
\renewcommand{\div}[1]{\vec{\nabla} \cdot{#1}}

\DeclareMathOperator{\sech}{sech}
\DeclareMathOperator{\cosech}{cosech}
\DeclareMathOperator{\arcsec}{arcsec}
\DeclareMathOperator{\arccot}{arcCot}
\DeclareMathOperator{\arccsc}{arcCsc}
\DeclareMathOperator{\arccosh}{arcCosh}
\DeclareMathOperator{\arcsinh}{arcsinh}
\DeclareMathOperator{\arctanh}{arctanh}
\DeclareMathOperator{\arcsech}{arcsech}
\DeclareMathOperator{\arccsch}{arcCsch}
\DeclareMathOperator{\arccoth}{arcCoth}
\DeclareMathOperator{\Tr}{Tr}

% Margin/inline comment macros used across the original manuscripts.
% Rendered in color in draft mode; harmless to leave in for the thesis.
\makeatletter
\DeclareRobustCommand{\rc}[1]{\textcolor{olive}{#1}}
\DeclareRobustCommand{\rg}[1]{\textcolor{magenta}{#1}}
\DeclareRobustCommand{\rgc}[1]{\textcolor{magenta}{{[\bf #1]}}}
\DeclareRobustCommand{\rcomm}[1]{\textcolor{OliveGreen}{[#1]}}
\DeclareRobustCommand{\revise}[1]{\textcolor{blue}{#1}}
\makeatother
\newcommand{\hjc}[1]{\textcolor{blue}{#1}}

% Chapter-specific macros that appear inside figure or table captions. They
% have to live here rather than in their chapter, because \listoffigures and
% \listoftables typeset the .lof and .lot files in the front matter, long
% before the chapters that define them are read.
% Internal Tide-Eddy run names:
\newcommand{\op}[1]{\operatorname{#1}}
\newcommand{\roeff}{Ro_{\op{eff}}}
\newcommand{\lr}{\texttt{Ro-L}\xspace}
\newcommand{\ir}{\texttt{Ro-I}\xspace}
\newcommand{\hr}{\texttt{Ro-H}\xspace}
\newcommand{\lre}{\texttt{Ro-L8}\xspace}
\newcommand{\lro}{\texttt{Ro-L1}\xspace}
\newcommand{\ire}{\texttt{Ro-I8}\xspace}
\newcommand{\iro}{\texttt{Ro-I1}\xspace}
\newcommand{\hre}{\texttt{Ro-H8}\xspace}
\newcommand{\hro}{\texttt{Ro-H1}\xspace}
\newcommand{\filler}[1]{{\centering\bfseries #1\par}}
% Inertial Caustics particle labels:
\newcommand{\cs}{C~}
\newcommand{\su}{S~}


% \title{PhD Thesis}
% \author{Rajarshi Chattopadhyay}
% \date{\today}x
\DeclareRobustCommand{\rg}[1]{{\textcolor{magenta}{#1}}}
\DeclareRobustCommand{\rgc}[1]{\textcolor{magenta}{{[ \textbf{#1}]}}}
\begin{document}

\frontmatter
% \maketitle



\begin{titlepage}
    \centering
    \thispagestyle{empty}
    
    \vspace*{1.2cm}
    
    % --- Title ---
    {\Huge\bfseries
    New Dynamics in Microscopic and Macroscopic Geophysical Flows\par}
    
    \vspace{3.5cm}
    
    % --- A Thesis ---
    {\large A Thesis\par}
    
    \vspace{2.5cm}
    
    % --- Submission block ---
    {\large
    Submitted to the\\[0.2em]
    Tata Institute of Fundamental Research, Mumbai\\[0.2em]
    Subject Board of Physics\\[0.2em]
    for the degree of Doctor of Philosophy\par}
    
    \vspace{1.5cm}
    
    {\large by\par}
    
    \vspace{1cm}
    
    {\large\bfseries Rajarshi Chattopadhyay\par}
    
    \vfill
    
    % --- Institute / date ---
    {\large
    International Centre for Theoretical Sciences,\\[0.2em]
    Bengaluru\\[0.2em]
    Tata Institute of Fundamental Research\\[0.2em]
    \today \par}
    
    \vspace{1.5cm}
    
  \end{titlepage}
  % \clearpage
  % \thispagestyle{empty}
  % \null
  % \clearpage
  \clearpage
  \thispagestyle{empty}
  \begin{center}
  {\LARGE\bfseries DECLARATION\par}
  \end{center}
  
  \vspace{2.5cm}
  
  \setlength{\parindent}{1.5em}
  
  \noindent This thesis is a presentation of my original research work. Wherever
  contributions of others are involved, every effort is made to indicate
  this clearly, with due reference to the literature, and acknowledgement
  of collaborative research and discussions.
  
  \vspace{1em}
  
  \noindent The work was done under the guidance of Professor Rama Govindarajan at
  the International Centre for Theoretical Sciences, Tata Institute of
  Fundamental Research, Bengaluru. Professor Jim Thomas served as the co-supervisor.
  
  \vspace{3cm}
  
  {\raggedleft
  \textbf{Rajarshi Chattopadhyay}\par
  }
  \hfill Date: \today
  
  \vspace{2.5cm}
  
  \noindent In my capacity as supervisor of the candidate's thesis, I certify that
  the above statements are true to the best of my knowledge.
  
  \vspace{3cm}
  
  \noindent\textbf{Professor Rama Govindarajan}
  
  \noindent Date: 
  
  \clearpage
  
    
  \newpage
  \thispagestyle{empty}   % keeps the roman page number "i" at the bottom
  \vspace*{0.35\textheight} % pushes text down to roughly the middle






\begingroup
\renewcommand{\chaptermark}[1]{}
\chapter*{Abstract}
This thesis studies turbulent flows and the particles suspended in them, focusing on problems where the enormous range of length scales involved makes a full description impossible. The transport of heat and mixing in the ocean is one motivation, where energy exchange between slowly moving eddies and fast internal waves spans scales from thousands of kilometers down to a few meters. We first ask how internal tides, exchange energy with ocean eddies across length scales, and find that at scales of a few kilometers this exchange is strong enough to spread energy across a broad range of frequencies, forming a continuum. This is absent at scales of hundreds of kilometers, pointing to an energy pathway from internal tides to dissipation that current ocean models do not resolve. We then study the Galerkin-truncated Euler equation, obtained by setting viscosity to zero and keeping only a finite number of modes, which thermalizes into an equilibrium state. The time taken to thermalize does not depend on the number of modes kept, and the energy flux across scales stays constant even without a Kolmogorov $k^{-5/3}$ spectrum, a result tied directly to the finite truncation.

The second motivation is the growth of cloud droplets into raindrops, where collisions between droplets of few microns impact kilometer-sized clouds. We study collisions, or caustics, of particles denser than the surrounding fluid. We find that although high strain regions promote caustics, there are that survive caustics even after encountering strong strain. The differnce in caustics and survivor particles is the type of strain they experience earlier in their history, and denser particles need less strain to form caustics. In a collaborative work, we then ask whether self-propelled, or active, particles collide for similar reasons despite having no inertia, and find that their self-propulsion acts like an effective mass, making their collisions similar to inertial particles. Finally, moving beyond particles of a single size, we build a combined Euler-Lagrange model for the growth of water droplets of different sizes in turbulence, where a few large droplets that detach from the flow sweep up a continuum of smaller ones. The model shows that smaller lucky droplets grow faster relative to their own size in 2D, but behaviour flips in 3D.

% TODO: thesis-level abstract summarizing all chapters.
\endgroup

\tableofcontents

% Both lists are rebuilt from the .lof and .lot files on every run, so adding
% or removing a figure or table updates them on the next compile.
\cleardoublepage
\phantomsection
\addcontentsline{toc}{chapter}{\listfigurename}
\listoffigures

\cleardoublepage
\phantomsection
\addcontentsline{toc}{chapter}{\listtablename}
\listoftables

\cleardoublepage
\phantomsection
\addcontentsline{toc}{chapter}{List of Symbols}
% % List of symbols, long version. Every symbol a reader of this thesis could need.
%
% THERE ARE TWO LISTS. The short version is symbols_short.tex, and it is a strict
% subset of this file with identical wording. A correction made here has to be
% made there as well, and the other way round.
%
% Loaded from main.tex between \listoftables and \mainmatter. Needs longtable and
% booktabs, both loaded in preamble.tex.
%
% Symbols are written in the form the reader sees, not the form the chapter
% source writes. The Tides chapter redefines \hat as \widehat, \bar as \overline
% and \t as \widetilde inside its own group, so those accents are spelt out here.
%
% The chapter column prints a chapter number through \ref. The labels are
%   \ref{tide-eddy}    Internal tide-eddy interactions
%   \ref{trunc-euler}  Thermalisation of the truncated Euler equation
%   \ref{fn-cstcs}     Caustics of finitely dense inertial particles
%   \ref{act-cstcs}    Active caustics
%   \ref{lck-drp}      Lucky drops
% A row reading "all" belongs to every chapter.

\chapter*{List of Symbols}
\markboth{List of Symbols}{List of Symbols}

\begingroup
\setlength{\LTpre}{0pt}
\setlength{\LTpost}{0pt}
\renewcommand{\arraystretch}{1.15}

\begin{longtable}{@{}p{0.18\textwidth}p{0.62\textwidth}p{0.12\textwidth}@{}}
\toprule
\textbf{Symbol} & \textbf{Meaning} & \textbf{Chapter} \\
\midrule
\endfirsthead
\toprule
\textbf{Symbol} & \textbf{Meaning} & \textbf{Chapter} \\
\midrule
\endhead
\midrule
\multicolumn{3}{r@{}}{\textit{continued on the next page}} \\
\endfoot
\bottomrule
\endlastfoot

\multicolumn{3}{@{}l}{\textbf{Latin letters}} \\
\addlinespace
$\mathbb{A}$ & fluid velocity gradient tensor, $St\,\nabla\boldsymbol{u}$ & \ref{fn-cstcs} \\
$\mathbf{A}$, $\tilde{\mathbf{A}}$ & antisymmetric part of the velocity gradient, and its non-dimensional form $\tau_{\rm act}\mathbf{A}$ & \ref{act-cstcs} \\
$A$ & avoider particles, which never meet high strain & \ref{fn-cstcs} \\
$A_a$, $A_g$ & amplitude of the tidal and the geostrophic forcing & \ref{tide-eddy} \\
$\AR$ & ratio of buoyancy to inertial frequency, $N/f$, set to $20$ & \ref{tide-eddy} \\
$a$ & radius of a particle & \ref{fn-cstcs} \\
$a$ & first invariant of $\mathbb{A}$, equal to $\Tr\mathbb{A} = 0$ & \ref{fn-cstcs} \\
$a_0$, $a_1$, $a_2$ & coefficients of the quartic in $\chi$ & \ref{fn-cstcs} \\
$a_s$, $a_b$ & radius of a small and of a big droplet & \ref{lck-drp} \\
\addlinespace
$\boldsymbol{B}$, $B_k$ & subfilter buoyancy flux & \ref{tide-eddy} \\
$b$ & buoyancy & \ref{tide-eddy} \\
\addlinespace
$C$ & caustics particles, which collide & \ref{fn-cstcs} \\
$\mathcal{C}$ & time-dependent constant in $\mathbb{E}(k) = \mathcal{C}k^2$ & \ref{trunc-euler} \\
$\mathcal{C}_\infty$ & its value once the flow has fully thermalised & \ref{trunc-euler} \\
$c$, $c_s$, $c_b$ & volume fraction of droplets, small ones, big ones & \ref{lck-drp} \\
\addlinespace
$D$, $\widehat{D}(k_h,t)$ & net flow dissipation and its spectrum & \ref{tide-eddy} \\
$D_\eta$ & noise strength in $\sqrt{2D_\eta}\,\boldsymbol{\eta}$ & \ref{act-cstcs} \\
$D/Dt$ & material derivative following the fluid & \ref{fn-cstcs}, \ref{lck-drp} \\
$D_t$ & fluid material derivative, $\partial_t + \boldsymbol{u}\cdot\nabla$ & \ref{act-cstcs} \\
$d$ & particle length scale used to non-dimensionalise & \ref{act-cstcs} \\
$\mathfrak{d}$ & $\det(\mathbb{Z})$, an invariant of the particle gradient & \ref{fn-cstcs} \\
$d_t$ & particle-phase material derivative, $\partial_t + v_k\partial_k$ & \ref{act-cstcs} \\
$d_s$ & diameter of a small droplet & \ref{lck-drp} \\
\addlinespace
$E$ & domain-integrated energy & \ref{tide-eddy} \\
$\widehat{E}(k_h,t)$ & horizontal energy spectrum & \ref{tide-eddy} \\
$E_g$, $E_a$ & geostrophic and ageostrophic energy & \ref{tide-eddy} \\
$E_s$, $E_f$ & slow and fast energy & \ref{tide-eddy} \\
$E_c^g$, $E_c^a$ & energy at which each forcing switches off & \ref{tide-eddy} \\
$E_{lw}$ & energy held by the linear waves & \ref{tide-eddy} \\
$\mathbb{E}$, $\mathbb{E}_0$ & kinetic energy per unit mass, and its initial value & \ref{trunc-euler} \\
$\mathbb{E}(k)$ & isotropic energy spectrum & \ref{trunc-euler} \\
$\widehat{e}$ & flow energy at one wavevector & \ref{tide-eddy} \\
$\widetilde{e}$ & pointwise energy of the filtered large-scale flow & \ref{tide-eddy} \\
$\boldsymbol{e}_\pm$ & unit vectors along the extensional and compressional strain & \ref{fn-cstcs} \\
$\boldsymbol{e}_r$ & radial unit vector & \ref{act-cstcs}, \ref{lck-drp} \\
$\boldsymbol{e}_\vartheta$ & azimuthal unit vector & \ref{fn-cstcs}, \ref{act-cstcs}, \ref{lck-drp} \\
$\boldsymbol{e}_v$ & unit vector along the particle velocity & \ref{fn-cstcs} \\
\addlinespace
$\boldsymbol{F}_{\boldsymbol{v}_{\rm 2D}}$, $F_w$, $F_b$ & forcing terms of the Boussinesq equations & \ref{tide-eddy} \\
$\widehat{F}(k_h,t)$ & power input spectrum & \ref{tide-eddy} \\
$\boldsymbol{F}_{v_g}$, $F_{w_g}$, $F_{b_g}$ & geostrophic forcing applied, the pattern times the amplitude $A_g$ times a switch & \ref{tide-eddy} \\
$\widehat{F}_{u_a}$, $\widehat{F}_{v_a}$, $\widehat{F}_{w_a}$, $\widehat{F}_{b_a}$ & ageostrophic forcing in spectral space & \ref{tide-eddy} \\
$F_\zeta$ & vorticity forcing of the two-dimensional flow & \ref{fn-cstcs} \\
$F_0$ & amplitude of the deterministic forcing on the background flow & \ref{act-cstcs} \\
$\boldsymbol{F}_{\rm eff}$ & effective external force on the particle & \ref{act-cstcs} \\
$\boldsymbol{F}_{\rm ext}$ & external force on the particle & \ref{act-cstcs} \\
$F_h$ & fluid drag on a big droplet & \ref{lck-drp} \\
$f$ & inertial, or Coriolis, frequency & \ref{tide-eddy} \\
$f_{p_g}$, $\boldsymbol{f}_{v_g}$, $f_{w_g}$, $f_{b_g}$ & geostrophic forcing patterns, functions of position alone, with $f_{w_g} = 0$ & \ref{tide-eddy} \\
$\widehat{f}_{p_a}$ & tidal forcing pressure in spectral space & \ref{tide-eddy} \\
$\widehat{f}_{u_a}$, $\widehat{f}_{v_a}$, $\widehat{f}_{w_a}$, $\widehat{f}_{b_a}$ & tidal forcing patterns in spectral space & \ref{tide-eddy} \\
\addlinespace
$g$ & logarithm of the small-droplet number density & \ref{lck-drp} \\
\addlinespace
$\mathbb{H}$ & Hamiltonian of the Euler equation & \ref{trunc-euler} \\
$\mathbb{H}$ & effective Hamiltonian of the inner solution & \ref{act-cstcs} \\
$\mathcal{H}$ & helicity & \ref{trunc-euler} \\
$H$ & vertical length scale & \ref{tide-eddy} \\
$h$ & H\"older regularity exponent of the velocity & \ref{trunc-euler} \\
$h(s)$ & function whose root fixes the clipping shift & \ref{lck-drp} \\
\addlinespace
$\mathbb{I}$ & identity tensor & \ref{act-cstcs} \\
$\mathbb{I}_4$ & four-dimensional identity matrix & \ref{tide-eddy} \\
$i$ & the imaginary unit, $\sqrt{-1}$ & \ref{trunc-euler} \\
$i$, $j$ & cell indices of the density-fluctuation measure & \ref{act-cstcs} \\
\addlinespace
$j^{\rm in}$, $j^{\rm out}$ & influx and outflux of droplet probability & \ref{lck-drp} \\
\addlinespace
$\mathbf{K}_{\rm G}$, $K_{\rm G}$ & Galerkin wavevector and truncation wavenumber & \ref{trunc-euler} \\
$K_{\rm I}$ & wavenumber where $k^{-5/3}$ gives way to $k^{-3}$ & \ref{trunc-euler} \\
$K_{\rm th}$ & wavenumber beyond which the spectrum has thermalised & \ref{trunc-euler} \\
$\mathbf{k}$, $k$ & wavevector and its magnitude & \ref{trunc-euler} \\
$k$ & total wavenumber, $\sqrt{k_h^2+k_z^2}$ & \ref{tide-eddy} \\
$k_0$ & wavenumber holding the energy at the start & \ref{trunc-euler} \\
$\boldsymbol{k}_h$, $k_h$ & horizontal wavevector $(k_x,k_y)$ and its magnitude & \ref{tide-eddy} \\
$k_z$ & vertical wavenumber & \ref{tide-eddy} \\
$k_\ell$ & filtering wavenumber, $2\pi/\ell$ & \ref{tide-eddy} \\
$k_m$ & largest wavenumber present in the system & \ref{tide-eddy} \\
$k_f$ & forcing wavenumber of the two-dimensional flow & \ref{fn-cstcs}, \ref{act-cstcs} \\
$k_{\mathrm{m}}$ & largest de-aliased wavenumber of the simulation & \ref{lck-drp} \\
\addlinespace
$\mathscr{L}$ & angular momentum per unit mass & \ref{fn-cstcs}, \ref{act-cstcs} \\
$\mathscr{L}^*$ & angular momentum at the matching time $t^*$ & \ref{act-cstcs} \\
$\mathscr{L}_f$ & final angular momentum of the outer solution, equal to $1$ & \ref{act-cstcs} \\
$L$ & horizontal length scale & \ref{tide-eddy} \\
$L$ & flow length scale & \ref{fn-cstcs}, \ref{lck-drp} \\
$L$ & length scale $\sqrt{\tilde{\Gamma}\tau_{\rm act}}$ of the point-vortex problem & \ref{act-cstcs} \\
$L_{box}$ & box size in dimensional units & \ref{lck-drp} \\
$l$ & polar flow-alignment length of the dimer & \ref{act-cstcs} \\
$\mathfrak{l}$ & size of the inner region, $\mathfrak{l}^2 = \tau\sqrt{3\alpha-2}$ & \ref{fn-cstcs} \\
$\ell$ & filtering length scale & \ref{tide-eddy} \\
\addlinespace
$M$ & droplet mass a reader asks the growth to reach & \ref{lck-drp} \\
$M_b$, $M_s$, $M_T$ & mass of all big drops, all small ones, and both & \ref{lck-drp} \\
$m$ & droplet mass & \ref{lck-drp} \\
$\mathrm{m}$ & parameter setting the Kolmogorov length, $\eta_s = \mathrm{m}/k_{\mathrm{m}}$ & \ref{lck-drp} \\
$m_b$, $m_p$ & mass of a big droplet and of any particle & \ref{lck-drp} \\
\addlinespace
$N$ & buoyancy frequency & \ref{tide-eddy} \\
$N_{\rm bin}$ & cells per direction in the density-fluctuation measure & \ref{act-cstcs} \\
$N_{\rm G}$ & number of retained modes & \ref{trunc-euler} \\
$N_g$ & collocation points per direction & \ref{trunc-euler} \\
$N_g$ & grid resolution & \ref{lck-drp} \\
$N_p$ & number of particles & \ref{act-cstcs} \\
$N_b$ & number of lucky, that is big, drops & \ref{lck-drp} \\
$\dot{N}_c$, $\dot{N}_{cl}$ & collision and coalescence rate of a big droplet & \ref{lck-drp} \\
$N_s(0)$ & number of small droplets at the start & \ref{lck-drp} \\
$N_V$ & number of particles per unit volume & \ref{lck-drp} \\
$N_t$ & number of time samples in the window & \ref{tide-eddy} \\
$n$ & particle number density & \ref{fn-cstcs}, \ref{lck-drp} \\
$n_s$, $\tilde{n}_s$ & small-droplet number density, and its scaled form in the code & \ref{lck-drp} \\
\addlinespace
$\tilde{\mathbf{O}}$ & fluid velocity gradient tensor, $\tau_{\rm act}\nabla\boldsymbol{u}$ & \ref{act-cstcs} \\
\addlinespace
$\mathbb{P}$ & projection operator onto the geostrophic flow & \ref{tide-eddy} \\
$\mathbb{P}_{\mathbf{K}_{\rm G}}$ & Galerkin projector & \ref{trunc-euler} \\
$\mathcal{P}$ & pressure & \ref{tide-eddy}, \ref{trunc-euler} \\
$\widehat{\mathcal{P}}_{\mathbf{k}}$ & Fourier mode of the pressure & \ref{trunc-euler} \\
$P(m,t)$ & probability density of a droplet of mass $m$ & \ref{lck-drp} \\
$P_{cc}$ & probability that a collision ends in coalescence & \ref{lck-drp} \\
$P_{\boldsymbol{x}}(n)$ & Poisson law for the droplet count in a cell & \ref{lck-drp} \\
$p_r$ & radial momentum per unit mass, $d\tilde{r}/d\tilde{t}$ & \ref{act-cstcs} \\
$\mathbf{p}$, $\mathbf{q}$ & wavevectors in the convolution sum & \ref{trunc-euler} \\
$\boldsymbol{p}$ & linear momentum & \ref{lck-drp} \\
\addlinespace
$Q$ & Okubo-Weiss parameter, $-\Tr(\mathbb{A}^2)/2$ & \ref{fn-cstcs}, \ref{lck-drp} \\
$Q_m$ & lowest $Q$ met by one caustics particle & \ref{fn-cstcs} \\
$\langle Q\rangle_{min}$ & lowest trajectory-averaged $Q$ before caustics & \ref{fn-cstcs} \\
$\langle Q_m\rangle_c$ & mean of $Q_m$ over the caustics particles & \ref{fn-cstcs} \\
$Q_{\rm rms}$ & root mean square of the Okubo-Weiss parameter & \ref{lck-drp} \\
\addlinespace
$R$ & third invariant of the fluid gradient, $-\det\mathbb{A}$ & \ref{fn-cstcs} \\
$R_b$ & radii of all the lucky drops, added up & \ref{lck-drp} \\
$\widehat{\boldsymbol{R}}_{\boldsymbol{v}_{\rm 2D}}$, $\widehat{R}_w$, $\widehat{R}_b$ & Fourier transforms of the nonlinear terms & \ref{tide-eddy} \\
$\mathcal{R}$, $\mathcal{R}_0$ & radius in inner units, $r/\mathfrak{l}$, and its initial value & \ref{fn-cstcs} \\
$\mathcal{R}_c$ & caustics radius in inner units & \ref{fn-cstcs} \\
$\mathcal{R}_{c0}$ & critical initial radius beyond which caustics never form & \ref{fn-cstcs}, \ref{act-cstcs} \\
$\operatorname{Re}$ & the real part & \ref{tide-eddy} \\
$Re$ & Reynolds number & \ref{act-cstcs}, \ref{lck-drp} \\
$Re_p$ & particle Reynolds number & \ref{fn-cstcs} \\
$Ro$, $Ro_{\rm eff}$ & Rossby number and effective Rossby number & \ref{tide-eddy} \\
$r$, $r_0$ & radial distance from the vortex, and its initial value & \ref{fn-cstcs}, \ref{act-cstcs}, \ref{lck-drp} \\
$r^*$ & radial distance at the matching time $t^*$ & \ref{act-cstcs} \\
$r_c$ & caustics radius & \ref{fn-cstcs} \\
$r_f$ & radius the particle settles at & \ref{act-cstcs} \\
\addlinespace
$S$ & survivor particles, which meet high strain yet avoid caustics & \ref{fn-cstcs} \\
$\mathbb{S}$ & strain tensor, $(\mathbb{A}+\mathbb{A}^{\!\top})/2$ & \ref{fn-cstcs} \\
$\mathbf{S}$, $\tilde{\mathbf{S}}$ & strain-rate tensor and its non-dimensional form $\tau_{\rm act}\mathbf{S}$ & \ref{act-cstcs} \\
$S(m,t)$ & probability that a droplet of mass $m$ survives & \ref{lck-drp} \\
$St$ & Stokes number & \ref{fn-cstcs}, \ref{act-cstcs}, \ref{lck-drp} \\
$St_s$, $St_b$ & Stokes number of the small and the big droplets & \ref{lck-drp} \\
$s^*$ & shift that restores the mean after clipping & \ref{lck-drp} \\
\addlinespace
$T_i$ & nonlinear energy transfer of flow component $i$ & \ref{tide-eddy} \\
$\widehat{T}(k_h,t)$ & nonlinear transfer spectrum & \ref{tide-eddy} \\
$T(k)$ & nonlinear transfer term & \ref{trunc-euler} \\
$T$ & timescale $\tau_{\rm act}$ of the point-vortex problem & \ref{act-cstcs} \\
$T$ & dimensional timescale of the simulation & \ref{lck-drp} \\
$\mathcal{T}$, $\mathcal{T}_c$ & time in inner units, $t/\tau$, and the caustics time in them & \ref{fn-cstcs} \\
$t$ & time & all \\
$t_i$, $t_f$ & start and end of the averaging window & \ref{tide-eddy} \\
$t_b$ & birth time of the oscillatory structures & \ref{trunc-euler} \\
$t_{\rm cas}$ & time the cascade takes to complete & \ref{trunc-euler} \\
$t_{\rm th}$ & time at which thermalisation sets in & \ref{trunc-euler} \\
$t_*$ & preshock time of the inviscid Burgers equation & \ref{trunc-euler} \\
$t_c$ & caustics time & \ref{fn-cstcs} \\
$t_{\rm thresh}$ & time at which $-\mathbb{Z}^2$ starts to dominate & \ref{fn-cstcs} \\
$t_m$, $t_{min}$ & time of the lowest $Q$, for one particle and on average & \ref{fn-cstcs} \\
$t^*$ & matching time for the outer boundary conditions & \ref{act-cstcs} \\
$t_x$ & time at which two adjacent rays cross & \ref{act-cstcs} \\
\addlinespace
$U$ & horizontal velocity scale & \ref{tide-eddy} \\
$U_0$ & scale of the background flow & \ref{act-cstcs} \\
$\boldsymbol{u}$ & three-dimensional fluid velocity $(u,v,w)$ & \ref{tide-eddy} \\
$\boldsymbol{u}$ & untruncated Euler velocity & \ref{trunc-euler} \\
$\mathbf{u}$, $\mathbf{u}_0$ & Galerkin-truncated velocity, and its initial value & \ref{trunc-euler} \\
$\widehat{\mathbf{u}}_{\mathbf{k}}$ & Fourier mode of the truncated velocity & \ref{trunc-euler} \\
$\boldsymbol{u}$ & background fluid velocity at the particle & \ref{fn-cstcs}, \ref{act-cstcs}, \ref{lck-drp} \\
$\hat{u}_k$ & scalar amplitude of the velocity at the wavenumber $k$ & \ref{trunc-euler} \\
$u$, $v$ & horizontal velocity components & \ref{tide-eddy} \\
$u_i$ & component of the velocity in index notation & \ref{tide-eddy} \\
$u_x$, $u_y$ & components of the two-dimensional flow & \ref{fn-cstcs} \\
$u_{\rm rms}$ & root mean square flow velocity & \ref{lck-drp} \\
\addlinespace
$V$ & domain of integration & \ref{tide-eddy} \\
$V_b$, $V_s$ & volume of the big and of the small droplets & \ref{lck-drp} \\
$\boldsymbol{v}$ & particle velocity & \ref{fn-cstcs}, \ref{act-cstcs}, \ref{lck-drp} \\
$\boldsymbol{v}_{\rm 2D}$ & horizontal fluid velocity, $(u,v)$ & \ref{tide-eddy} \\
$\boldsymbol{v}_s$, $\boldsymbol{v}_b$ & velocity of the small and of the big droplets & \ref{lck-drp} \\
$v_r$, $v_r^*$ & radial particle velocity, and its value at the matching time & \ref{act-cstcs} \\
\addlinespace
$\mathcal{W}$ & Okubo-Weiss parameter from dimensional gradients & \ref{act-cstcs} \\
$\mathbf{w}$ & end-to-end extension of the dimer & \ref{act-cstcs} \\
$\mathrm{w}_0$ & preferred length of the dimer & \ref{act-cstcs} \\
$w$ & vertical velocity & \ref{tide-eddy} \\
\addlinespace
$\mathbf{X}$ & centroid of the active dimer & \ref{act-cstcs} \\
$\mathcal{X}$ & percentage of $K_{\rm G}$ setting the fitting window & \ref{trunc-euler} \\
$\mathbf{x}$ & position & \ref{trunc-euler} \\
$\boldsymbol{x}$ & position & \ref{lck-drp} \\
$\boldsymbol{x}_p$ & position of a particle & \ref{fn-cstcs} \\
$\boldsymbol{x}_b$ & position of a big droplet & \ref{lck-drp} \\
$\widehat{\boldsymbol{x}}$, $\widehat{\boldsymbol{y}}$ & horizontal unit vectors & \ref{tide-eddy} \\
$x$, $y$, $z$ & spatial coordinates & all \\
\addlinespace
$\mathbb{Z}$ & particle velocity gradient tensor, $St\,\nabla\boldsymbol{v}$ & \ref{fn-cstcs} \\
$\mathbb{Z}_0$ & fixed point of the frozen-particle equation & \ref{fn-cstcs} \\
$\mathbb{Z}_{\rm thresh}$ & that tensor at the time $t_{\rm thresh}$ & \ref{fn-cstcs} \\
$\mathbf{Z}$ & particle velocity gradient tensor, $\tau_{\rm act}\nabla\boldsymbol{v}$ & \ref{act-cstcs} \\
$\mathcal{Z}$, $\mathcal{Z}_\infty$ & enstrophy, and the value it saturates at & \ref{trunc-euler} \\
$\widehat{\boldsymbol{z}}$ & vertical unit vector & \ref{tide-eddy} \\
$z_2$ & $\Tr(\mathbb{Z}^2)$, an invariant of the particle gradient & \ref{fn-cstcs} \\
\addlinespace
\midrule
\multicolumn{3}{@{}l}{\textbf{Greek letters}} \\
\addlinespace
$\alpha$ & density parameter, $2\rho/(2\rho+1)$ & \ref{fn-cstcs} \\
$\alpha_{\rm act}$ & Jeffery-type response of the dimer to the flow & \ref{act-cstcs} \\
$\alpha_g$ & fitted growth exponent in $t^{\alpha_g}$ & \ref{lck-drp} \\
\addlinespace
$\beta$ & swimming rate, so $1/\beta$ is the time to swim a body length & \ref{act-cstcs} \\
$\tilde{\beta}$ & non-dimensional motility, $\beta\sqrt{\tau_{\rm act}/\tilde{\Gamma}}$ & \ref{act-cstcs} \\
\addlinespace
$\Gamma$ & circulation of the point vortex & \ref{fn-cstcs}, \ref{lck-drp} \\
$\tilde{\Gamma}$ & circulation of the point vortex, $\Gamma/2\pi$ & \ref{act-cstcs} \\
$\gamma$ & spectral slope, $\mathbb{E}(k)\sim k^\gamma$ & \ref{trunc-euler} \\
$\gamma_\epsilon$ & linear drag coefficient of the two-dimensional flow & \ref{fn-cstcs} \\
\addlinespace
$\Delta$ & Laplacian operator & \ref{tide-eddy} \\
$\Delta\omega$ & width of the window around a dispersion curve & \ref{tide-eddy} \\
$\Delta_0\mathbb{E}$, $\Delta_0\mathcal{H}$ & relative errors in energy and helicity & \ref{trunc-euler} \\
$\Delta r$ & separation between two neighbouring particle rings & \ref{act-cstcs} \\
$\Delta t$ & numerical time step & \ref{trunc-euler}, \ref{act-cstcs}, \ref{lck-drp} \\
$\Delta\rho_{ij}$ & fluctuation of the number density in a cell & \ref{act-cstcs} \\
$\delta$ & trace of $\mathbb{Z}$, equal to $St\,\nabla\cdot\boldsymbol{v}$ & \ref{fn-cstcs} \\
$\delta(x)$ & Dirac delta function & \ref{lck-drp} \\
$\delta_{ij}$ & Kronecker delta & \ref{tide-eddy} \\
$\delta_i$, $\delta_o$ & spatial stretching factors of the inner and outer solutions & \ref{act-cstcs} \\
$\delta t$ & sampling interval of time in the discrete Fourier transform & \ref{tide-eddy} \\
$\delta\omega$ & frequency sampling interval & \ref{tide-eddy} \\
\addlinespace
$\epsilon$ & effective dissipation per unit mass & \ref{trunc-euler} \\
$\epsilon$ & mean energy dissipation rate of the cloud & \ref{lck-drp} \\
$\epsilon_{num}$ & dissipation rate reached by the simulation & \ref{lck-drp} \\
$\epsilon_{ijk}$ & Levi-Civita tensor & \ref{tide-eddy} \\
$\varepsilon$ & separation between two particle rings & \ref{fn-cstcs} \\
$\varepsilon_i$, $\varepsilon_o$ & temporal stretching factors of the inner and outer solutions & \ref{act-cstcs} \\
\addlinespace
$\zeta$ & vertical vorticity & \ref{tide-eddy}, \ref{fn-cstcs}, \ref{act-cstcs}, \ref{lck-drp} \\
$\zeta_{\rm rms}$ & root mean square vorticity & \ref{tide-eddy}, \ref{lck-drp} \\
\addlinespace
$\eta$, $\eta_s$ & Kolmogorov length, and its value in simulation units & \ref{lck-drp} \\
$\eta_{num}$ & Kolmogorov length set by the numerical viscosity & \ref{lck-drp} \\
$\boldsymbol{\eta}$ & white noise vector of the dimer equations & \ref{act-cstcs} \\
\addlinespace
$\Theta$ & Heaviside step function, switching the forcing off & \ref{tide-eddy} \\
$\theta$ & latitude & \ref{tide-eddy} \\
$\theta_{k_h}$ & azimuthal angle in horizontal wavenumber space & \ref{tide-eddy} \\
$\vartheta$ & polar angle of the point-vortex coordinates & \ref{fn-cstcs}, \ref{act-cstcs}, \ref{lck-drp} \\
$\dot{\vartheta}$ & angular velocity of the dimer, $d\vartheta/dt$ & \ref{act-cstcs} \\
$\vartheta_v$ & angle between particle velocity and the strain eigendirections & \ref{fn-cstcs} \\
\addlinespace
$\kappa$ & root mean square velocity gradient of the background flow & \ref{act-cstcs} \\
\addlinespace
$\varLambda(m,t)$ & total collision rate of a droplet of mass $m$ & \ref{lck-drp} \\
$\lambda_s$ & depletion rate of the small droplets & \ref{lck-drp} \\
$\lambda(m,m',t)$ & Smoluchowski collision kernel & \ref{lck-drp} \\
$\lambda(\boldsymbol{x})$ & local Poisson intensity set by the flow gradients & \ref{lck-drp} \\
$\bar{\lambda}$ & mean of the Poisson law, particles per grid cell & \ref{lck-drp} \\
\addlinespace
$\mu$ & Stokesian mobility of the particle & \ref{act-cstcs} \\
$\mu_E$ & Ekman friction of the background flow & \ref{act-cstcs} \\
$\mu$m & micrometre, a unit and not a symbol & \ref{act-cstcs}, \ref{lck-drp} \\
\addlinespace
$\nu$ & kinematic viscosity & \ref{trunc-euler}, \ref{fn-cstcs}, \ref{act-cstcs}, \ref{lck-drp} \\
$\nu_{\rm eff}$ & eddy viscosity & \ref{trunc-euler} \\
$\nu_{\rm hyp}$ & hyperviscosity multiplying $\Delta^8$ & \ref{tide-eddy} \\
$\nu_{num}$ & viscosity chosen for the simulation & \ref{lck-drp} \\
\addlinespace
$\varXi$, $\varXi_0$ & placeholder field whose flux fraction is computed, and its threshold & \ref{tide-eddy} \\
$\xi$ & shorthand for $1-\alpha$ & \ref{fn-cstcs} \\
\addlinespace
$\Pi$, $\Pi(k)$ & kinetic energy flux & \ref{trunc-euler} \\
$\widehat{\Pi}(k_h,t)$ & horizontal spectral flux & \ref{tide-eddy} \\
$\Pi_\ell$ & energy flux in physical space at the scale $\ell$ & \ref{tide-eddy} \\
\addlinespace
$\rho$ & particle-to-fluid density ratio, $\rho_p/\rho_f$ & \ref{fn-cstcs}, \ref{lck-drp} \\
$\rho_p$, $\rho_f$ & density of the particle and of the fluid & \ref{fn-cstcs}, \ref{lck-drp} \\
$\rho_{ij}$ & number of particles in the cell $(i,j)$ & \ref{act-cstcs} \\
\addlinespace
$\Sigma_\ell$ & sum of the divergence terms of the filtered budget & \ref{tide-eddy} \\
$\sigma$ & net strain, $\sqrt{\sigma_s^2+\sigma_n^2}$ & \ref{fn-cstcs} \\
$\sigma_n$, $\sigma_s$ & normal and shear strain & \ref{fn-cstcs} \\
$\varsigma$, $\varsigma_2$, $\varsigma_3$ & collision cross-section, in two and in three dimensions & \ref{lck-drp} \\
$\varsigma_b$ & cross-section of a big drop on the small ones, $\pi(a_s+a_b)^2$ & \ref{lck-drp} \\
\addlinespace
$\tau$ & small parameter of the inner expansion & \ref{fn-cstcs} \\
$\tau_{\rm act}$, $\tilde{\tau}_{\rm act}$ & relaxation time of the dimer, and its flow-dependent form & \ref{act-cstcs} \\
$\tau_p$, $\tau_f$ & particle relaxation time and flow timescale & \ref{fn-cstcs}, \ref{lck-drp} \\
$\tau_s$, $\tau_b$ & relaxation time of the small and the big droplets & \ref{lck-drp} \\
$\tau_\eta$ & Kolmogorov timescale & \ref{fn-cstcs}, \ref{lck-drp} \\
$\tau_{\eta_{num}}$ & Kolmogorov timescale in simulation units & \ref{lck-drp} \\
$\tau_{\rm coll}$ & collision timescale & \ref{lck-drp} \\
$\tau_d$ & eddy-damping timescale & \ref{trunc-euler} \\
$\tau_{ij}$, $\uptau$, $\boldsymbol{\tau}$, $\boldsymbol{\tau}_h$ & subfilter stress & \ref{tide-eddy} \\
\addlinespace
$\varphi$ & cumulative flux fraction & \ref{tide-eddy} \\
$\phi$, $\psi$ & arbitrary flow variables & \ref{tide-eddy} \\
\addlinespace
$\chi$ & auxiliary root of the quartic & \ref{fn-cstcs} \\
\addlinespace
$\Omega$ & rotation rate of the Earth & \ref{tide-eddy} \\
$\omega$ & wave frequency & \ref{tide-eddy} \\
$\omega_0$ & tidal forcing frequency & \ref{tide-eddy} \\
$\omega_c$ & cut-off frequency parting slow from fast & \ref{tide-eddy} \\
$\omega_{\rm max}$ & highest frequency the discrete transform resolves & \ref{tide-eddy} \\
\addlinespace
\midrule
\multicolumn{3}{@{}l}{\textbf{Operators, accents and conventions}} \\
\addlinespace
$\nabla$ & horizontal gradient, $\widehat{\boldsymbol{x}}\partial_x + \widehat{\boldsymbol{y}}\partial_y$ & \ref{tide-eddy} \\
$\nabla_3$ & three-dimensional gradient & \ref{tide-eddy} \\
$\nabla$ & gradient & \ref{fn-cstcs}, \ref{act-cstcs}, \ref{lck-drp} \\
\addlinespace
$\widehat{\phantom{u}}\;$ & a Fourier amplitude, as in $\widehat{u}$ & \ref{tide-eddy} \\
$\hat{\phantom{u}}\;$ & a Fourier amplitude, as in $\widehat{\mathbf{u}}_{\mathbf{k}}$ & \ref{trunc-euler} \\
$\overline{\phantom{u}}\;$ & a time average & \ref{tide-eddy} \\
$\langle\cdot\rangle$ & an average & \ref{fn-cstcs}, \ref{act-cstcs}, \ref{lck-drp} \\
$\langle\cdot\rangle_{\rm rms}$ & a root mean square average & \ref{act-cstcs} \\
$\langle\cdot\rangle_x$ & a spatial average & \ref{lck-drp} \\
$\widetilde{\phantom{u}}\;$ & the large-scale part left by the spatial filter & \ref{tide-eddy} \\
$\tilde{\phantom{u}}\;$ & a rescaled or non-dimensional quantity & \ref{act-cstcs} \\
$(\,)'$ & the small-scale part left by the spatial filter & \ref{tide-eddy} \\
$\dot{(\,)}$ & a derivative with respect to time & \ref{fn-cstcs}, \ref{act-cstcs}, \ref{lck-drp} \\
$\boldsymbol{u}$, bold italic & a vector & all \\
$\mathbf{u}$, bold upright & the Galerkin-truncated field, against the italic untruncated one & \ref{trunc-euler} \\
$\mathbf{S}$, bold upright & a tensor & \ref{act-cstcs} \\
$\mathbb{A}$, blackboard bold & a tensor & \ref{fn-cstcs} \\
$\mathcal{R}$, calligraphic & a variable of the inner region & \ref{fn-cstcs} \\
$(\,)^*$ & the complex conjugate & \ref{tide-eddy} \\
$(\,)^\bigstar$ & a dimensional quantity, as in $\omega^\bigstar$ & \ref{tide-eddy}, \ref{fn-cstcs} \\
$(\,)^\dagger$ & the Hermitian conjugate & \ref{tide-eddy} \\
$(\,)^\top$ & the transpose & \ref{fn-cstcs}, \ref{act-cstcs} \\
subscript $g$, $a$ & the geostrophic and the ageostrophic part & \ref{tide-eddy} \\
subscript $s$, $f$ & the slow and the fast part & \ref{tide-eddy} \\
subscript $\ell$ & taken at the filter scale $\ell$ & \ref{tide-eddy} \\
subscript $h$ & horizontal & \ref{tide-eddy} \\
subscript $p$, $f$ & particle and fluid & \ref{fn-cstcs} \\
subscript $s$, $b$ & the small and the big droplets & \ref{lck-drp} \\
subscript $0$ & a value at the initial time & \ref{fn-cstcs}, \ref{act-cstcs} \\
subscript $i$, $o$ & the inner and the outer solution & \ref{act-cstcs} \\
subscript rms & a root mean square value & \ref{tide-eddy}, \ref{act-cstcs}, \ref{lck-drp} \\
repeated index & summed over, by the Einstein convention & \ref{tide-eddy}, \ref{act-cstcs} \\

\end{longtable}
\endgroup

% List of symbols, short version. The principal symbols only.
%
% THERE ARE TWO LISTS. The long version is symbols_full.tex, and this file is a
% strict subset of it: every row below is copied from there word for word. A
% correction made here has to be made there as well, and the other way round.
%
% Loaded from main.tex between \listoftables and \mainmatter. Needs longtable and
% booktabs, both loaded in preamble.tex.
%
% Symbols are written in the form the reader sees, not the form the chapter
% source writes. The Tides chapter redefines \hat as \widehat, \bar as \overline
% and \t as \widetilde inside its own group, so those accents are spelt out here.
%
% The chapter column prints a chapter number through \ref. The labels are
%   \ref{tide-eddy}    Internal tide-eddy interactions
%   \ref{trunc-euler}  Thermalisation of the truncated Euler equation
%   \ref{fn-cstcs}     Caustics of finitely dense inertial particles
%   \ref{act-cstcs}    Active caustics
%   \ref{lck-drp}      Lucky drops
% A row reading "all" belongs to every chapter.

\chapter*{List of Symbols}
\markboth{List of Symbols}{List of Symbols}

\begingroup
\setlength{\LTpre}{0pt}
\setlength{\LTpost}{0pt}
\renewcommand{\arraystretch}{1.15}

\begin{longtable}{@{}p{0.18\textwidth}p{0.62\textwidth}p{0.12\textwidth}@{}}
\toprule
\textbf{Symbol} & \textbf{Meaning} & \textbf{Chapter} \\
\midrule
\endfirsthead
\toprule
\textbf{Symbol} & \textbf{Meaning} & \textbf{Chapter} \\
\midrule
\endhead
\midrule
\multicolumn{3}{r@{}}{\textit{continued on the next page}} \\
\endfoot
\bottomrule
\endlastfoot

\multicolumn{3}{@{}l}{\textbf{Latin letters}} \\
\addlinespace
$\mathbb{A}$ & fluid velocity gradient tensor, $St\,\nabla\boldsymbol{u}$ & \ref{fn-cstcs} \\
$\mathbf{A}$, $\tilde{\mathbf{A}}$ & antisymmetric part of the velocity gradient, and its non-dimensional form $\tau_{\rm act}\mathbf{A}$ & \ref{act-cstcs} \\
$A$ & avoider particles, which never meet high strain & \ref{fn-cstcs} \\
$\AR$ & ratio of buoyancy to inertial frequency, $N/f$, set to $20$ & \ref{tide-eddy} \\
$b$ & buoyancy & \ref{tide-eddy} \\
$C$ & caustics particles, which collide & \ref{fn-cstcs} \\
$c$, $c_s$, $c_b$ & volume fraction of droplets, small ones, big ones & \ref{lck-drp} \\
$D$, $\widehat{D}(k_h,t)$ & net flow dissipation and its spectrum & \ref{tide-eddy} \\
$D_\eta$ & noise strength in $\sqrt{2D_\eta}\,\boldsymbol{\eta}$ & \ref{act-cstcs} \\
$D/Dt, D_t$ & material derivative following the fluid  $\partial_t + \boldsymbol{u}\cdot\nabla$& \ref{fn-cstcs}, \ref{lck-drp}, \ref{act-cstcs} . \\
$E$ & domain-integrated energy & \ref{tide-eddy} \\
$\widehat{E}(k_h,t)$ & horizontal energy spectrum & \ref{tide-eddy} \\
$E_g$, $E_a$ & geostrophic and ageostrophic energy & \ref{tide-eddy} \\
$E_s$, $E_f$ & slow and fast energy & \ref{tide-eddy} \\
$\mathbb{E}$, $\mathbb{E}_0$ & kinetic energy per unit mass, and its initial value & \ref{trunc-euler} \\
$\mathbb{E}(k)$ & isotropic energy spectrum & \ref{trunc-euler} \\
$\boldsymbol{e}_\pm$ & unit vectors along the extensional and compressional strain & \ref{fn-cstcs} \\
$\boldsymbol{F}_{\boldsymbol{v}_{\rm 2D}}$, $F_w$, $F_b$ & forcing terms of the Boussinesq equations & \ref{tide-eddy} \\
$F_0$ & amplitude of the deterministic forcing on the background flow & \ref{act-cstcs} \\
$\boldsymbol{F}_{\rm ext}$ & external force on the particle & \ref{act-cstcs} \\
$f$ & inertial, or Coriolis, frequency & \ref{tide-eddy} \\
$\mathbb{H}$ & Hamiltonian of the Euler equation & \ref{trunc-euler} \\
$\mathbb{H}$ & effective Hamiltonian of the inner solution & \ref{act-cstcs} \\
$\mathcal{H}$ & helicity & \ref{trunc-euler} \\
$H$ & vertical length scale & \ref{tide-eddy} \\
$\mathbf{K}_{\rm G}$, $K_{\rm G}$ & Galerkin wavevector and truncation wavenumber & \ref{trunc-euler} \\
$\mathbf{k}$, $k$ & wavevector and its magnitude & \ref{trunc-euler} \\
$\boldsymbol{k}_h$, $k_h$ & horizontal wavevector $(k_x,k_y)$ and its magnitude & \ref{tide-eddy} \\
$k_z$ & vertical wavenumber & \ref{tide-eddy} \\
$k_f$ & forcing wavenumber of the two-dimensional flow & \ref{fn-cstcs}, \ref{act-cstcs} \\
$\mathscr{L}$ & angular momentum per unit mass & \ref{fn-cstcs}, \ref{act-cstcs} \\
$\mathscr{L}^*$ & angular momentum at the matching time $t^*$ & \ref{act-cstcs} \\
$\mathscr{L}_f$ & final angular momentum of the outer solution, equal to $1$ & \ref{act-cstcs} \\
$L$ & horizontal length scale & \ref{tide-eddy} \\
$L$ & flow length scale & \ref{fn-cstcs}, \ref{lck-drp} \\
$l$ & polar flow-alignment length of the dimer & \ref{act-cstcs} \\
$\mathfrak{l}$ & size of the inner region, $\mathfrak{l}^2 = \tau\sqrt{3\alpha-2}$ & \ref{fn-cstcs} \\
$\ell$ & filtering length scale & \ref{tide-eddy} \\
$m$ & droplet mass & \ref{lck-drp} \\
$\mathrm{m}$ & parameter setting the Kolmogorov length, $\eta_s = \mathrm{m}/k_{\mathrm{m}}$ & \ref{lck-drp} \\
$N$ & buoyancy frequency & \ref{tide-eddy} \\
$N_{\rm G}$ & number of retained modes & \ref{trunc-euler} \\
$N_g$ & collocation points per direction & \ref{trunc-euler} \\
$N_g$ & grid resolution & \ref{lck-drp} \\
$N_p$ & number of particles & \ref{act-cstcs} \\
$n$ & particle number density & \ref{fn-cstcs}, \ref{lck-drp} \\
$\tilde{\mathbf{O}}$ & fluid velocity gradient tensor, $\tau_{\rm act}\nabla\boldsymbol{u}$ & \ref{act-cstcs} \\
$\mathbb{P}_{\mathbf{K}_{\rm G}}$ & Galerkin projector & \ref{trunc-euler} \\
$\mathcal{P}$ & pressure & \ref{tide-eddy}, \ref{trunc-euler} \\
$\widehat{\mathcal{P}}_{\mathbf{k}}$ & Fourier mode of the pressure & \ref{trunc-euler} \\
$P(m,t)$ & probability density of a droplet of mass $m$ & \ref{lck-drp} \\
$p_r$ & radial momentum per unit mass, $d\tilde{r}/d\tilde{t}$ & \ref{act-cstcs} \\
$\mathbf{p}$, $\mathbf{q}$ & wavevectors in the convolution sum & \ref{trunc-euler} \\
$\boldsymbol{p}$ & linear momentum & \ref{lck-drp} \\
$Q$ & Okubo-Weiss parameter, $-\Tr(\mathbb{A}^2)/2$ & \ref{fn-cstcs}, \ref{lck-drp} \\
$R$ & third invariant of the fluid gradient, $-\det\mathbb{A}$ & \ref{fn-cstcs} \\
$\mathcal{R}$, $\mathcal{R}_0$ & radius in inner units, $r/\mathfrak{l}$, and its initial value & \ref{fn-cstcs} \\
$\mathcal{R}_{c0}$ & critical initial radius beyond which caustics never form & \ref{fn-cstcs}, \ref{act-cstcs} \\
$Re$ & Reynolds number & \ref{act-cstcs}, \ref{lck-drp} \\
$Ro$, $Ro_{\rm eff}$ & Rossby number and effective Rossby number & \ref{tide-eddy} \\
$r$, $r_0$ & radial distance from the vortex, and its initial value & \ref{fn-cstcs}, \ref{act-cstcs}, \ref{lck-drp} \\
$r^*$ & radial distance at the matching time $t^*$ & \ref{act-cstcs} \\
$r_c$ & caustics radius & \ref{fn-cstcs} \\
$S$ & survivor particles, which meet high strain yet avoid caustics & \ref{fn-cstcs} \\
$\mathbb{S}$ & strain tensor, $(\mathbb{A}+\mathbb{A}^{\!\top})/2$ & \ref{fn-cstcs} \\
$\mathbf{S}$, $\tilde{\mathbf{S}}$ & strain-rate tensor and its non-dimensional form $\tau_{\rm act}\mathbf{S}$ & \ref{act-cstcs} \\
$St$ & Stokes number & \ref{fn-cstcs}, \ref{act-cstcs}, \ref{lck-drp} \\
$St_s$, $St_b$ & Stokes number of the small and the big droplets & \ref{lck-drp} \\
$T_i$ & nonlinear energy transfer of flow component $i$ & \ref{tide-eddy} \\
$T(k)$ & nonlinear transfer term & \ref{trunc-euler} \\
$\mathcal{T}$, $\mathcal{T}_c$ & time in inner units, $t/\tau$, and the caustics time in them & \ref{fn-cstcs} \\
$t_b$ & birth time of the oscillatory structures & \ref{trunc-euler} \\
$t_{\rm cas}$ & time the cascade takes to complete & \ref{trunc-euler} \\
$t_{\rm th}$ & time at which thermalisation sets in & \ref{trunc-euler} \\
$t_c$ & caustics time & \ref{fn-cstcs} \\
$U$ & horizontal velocity scale & \ref{tide-eddy} \\
$\boldsymbol{u}$ & three-dimensional fluid velocity $(u,v,w)$ & \ref{tide-eddy},\ref{trunc-euler} \\
$\mathbf{u}$, $\mathbf{u}_0$ & Galerkin-truncated velocity, and its initial value & \ref{trunc-euler} \\
$\boldsymbol{u}$ & background fluid velocity at the particle & \ref{fn-cstcs}, \ref{act-cstcs}, \ref{lck-drp} \\
$\hat{u}_k$ & scalar amplitude of the velocity at the wavenumber $k$ & \ref{trunc-euler} \\
$u$, $v$ & horizontal velocity components & \ref{tide-eddy} \\
$u_x$, $u_y$ & components of the two-dimensional flow & \ref{fn-cstcs} \\
$V$ & domain of integration & \ref{tide-eddy} \\
$\boldsymbol{v}$ & particle velocity & \ref{fn-cstcs}, \ref{act-cstcs}, \ref{lck-drp} \\
$\boldsymbol{v}_{\rm 2D}$ & horizontal fluid velocity, $(u,v)$ & \ref{tide-eddy} \\
$\boldsymbol{v}_s$, $\boldsymbol{v}_b$ & velocity of the small and of the big droplets & \ref{lck-drp} \\
$\mathcal{W}$ & Okubo-Weiss parameter from dimensional gradients & \ref{act-cstcs} \\
$\mathbf{w}$ & end-to-end extension of the dimer & \ref{act-cstcs} \\
$\mathrm{w}_0$ & preferred length of the dimer & \ref{act-cstcs} \\
$w$ & vertical velocity & \ref{tide-eddy} \\
$\mathbf{X}$ & centroid of the active dimer & \ref{act-cstcs} \\
$\mathbb{Z}$ & particle velocity gradient tensor, $St\,\nabla\boldsymbol{v}$ & \ref{fn-cstcs} \\
$\mathbb{Z}_0$ & fixed point of the frozen-particle equation & \ref{fn-cstcs} \\
$\mathbf{Z}$ & particle velocity gradient tensor, $\tau_{\rm act}\nabla\boldsymbol{v}$ & \ref{act-cstcs} \\
$\mathcal{Z}$, $\mathcal{Z}_\infty$ & enstrophy, and the value it saturates at & \ref{trunc-euler} \\
\addlinespace
\midrule
\multicolumn{3}{@{}l}{\textbf{Greek letters}} \\
\addlinespace
$\alpha$ & density parameter, $2\rho/(2\rho+1)$ & \ref{fn-cstcs} \\
$\alpha_{\rm act}$ & Jeffery-type response of the dimer to the flow & \ref{act-cstcs} \\
$\beta$ & swimming rate, so $1/\beta$ is the time to swim a body length & \ref{act-cstcs} \\
$\tilde{\beta}$ & non-dimensional motility, $\beta\sqrt{\tau_{\rm act}/\tilde{\Gamma}}$ & \ref{act-cstcs} \\
$\Gamma$ & circulation of the point vortex & \ref{fn-cstcs}, \ref{lck-drp} \\
$\tilde{\Gamma}$ & circulation of the point vortex, $\Gamma/2\pi$ & \ref{act-cstcs} \\
$\gamma$ & spectral slope, $\mathbb{E}(k)\sim k^\gamma$ & \ref{trunc-euler} \\
$\gamma_\epsilon$ & linear drag coefficient of the two-dimensional flow & \ref{fn-cstcs} \\
$\Delta r$ & separation between two neighbouring particle rings & \ref{act-cstcs} \\
$\delta$ & trace of $\mathbb{Z}$, equal to $St\,\nabla\cdot\boldsymbol{v}$ & \ref{fn-cstcs} \\
$\delta(x)$ & Dirac delta function & \ref{lck-drp} \\
$\epsilon$ & effective dissipation per unit mass & \ref{trunc-euler} \\
$\epsilon$ & mean energy dissipation rate of the cloud & \ref{lck-drp} \\
$\varepsilon$ & separation between two particle rings & \ref{fn-cstcs} \\
$\zeta$ & vertical vorticity & \ref{tide-eddy}, \ref{fn-cstcs}, \ref{act-cstcs}, \ref{lck-drp} \\
$\boldsymbol{\eta}$ & white noise vector of the dimer equations & \ref{act-cstcs} \\
$\Theta$ & Heaviside step function, switching the forcing off & \ref{tide-eddy} \\
$\vartheta$ & polar angle of the point-vortex coordinates & \ref{fn-cstcs}, \ref{act-cstcs}, \ref{lck-drp} \\
$\dot{\vartheta}$ & angular velocity of the dimer, $d\vartheta/dt$ & \ref{act-cstcs} \\
$\kappa$ & root mean square velocity gradient of the background flow & \ref{act-cstcs} \\
$\lambda_s$ & depletion rate of the small droplets & \ref{lck-drp} \\
$\lambda(m,m',t)$ & Smoluchowski collision kernel & \ref{lck-drp} \\
$\mu$ & Stokesian mobility of the particle & \ref{act-cstcs} \\
$\mu_E$ & Ekman friction of the background flow & \ref{act-cstcs} \\
$\nu$ & kinematic viscosity & \ref{trunc-euler}, \ref{fn-cstcs}, \ref{act-cstcs}, \ref{lck-drp} \\
$\nu_{\rm hyp}$ & hyperviscosity multiplying $\Delta^8$ & \ref{tide-eddy} \\
$\varXi$, $\varXi_0$ & placeholder field whose flux fraction is computed, and its threshold & \ref{tide-eddy} \\
$\Pi$, $\Pi(k)$ & kinetic energy flux & \ref{trunc-euler} \\
$\widehat{\Pi}(k_h,t)$ & horizontal spectral flux & \ref{tide-eddy} \\
$\Pi_\ell$ & energy flux in physical space at the scale $\ell$ & \ref{tide-eddy} \\
$\rho$ & particle-to-fluid density ratio, $\rho_p/\rho_f$ & \ref{fn-cstcs}, \ref{lck-drp} \\
$\sigma$ & net strain, $\sqrt{\sigma_s^2+\sigma_n^2}$ & \ref{fn-cstcs} \\
$\varsigma$, $\varsigma_2$, $\varsigma_3$ & collision cross-section, in two and in three dimensions & \ref{lck-drp} \\
$\varsigma_b$ & cross-section of a big drop on the small ones, $\pi(a_s+a_b)^2$ & \ref{lck-drp} \\
$\tau$ & small parameter of the inner expansion & \ref{fn-cstcs} \\
$\tau_{\rm act}$, $\tilde{\tau}_{\rm act}$ & relaxation time of the dimer, and its flow-dependent form & \ref{act-cstcs} \\
$\tau_p$, $\tau_f$ & particle relaxation time and flow timescale & \ref{fn-cstcs}, \ref{lck-drp} \\
$\varphi$ & cumulative flux fraction & \ref{tide-eddy} \\
$\omega$ & wave frequency & \ref{tide-eddy} \\
$\omega_0$ & tidal forcing frequency & \ref{tide-eddy} \\
\addlinespace
\midrule
\multicolumn{3}{@{}l}{\textbf{Operators, accents and conventions}} \\
\addlinespace
$\nabla$ & horizontal gradient, $\widehat{\boldsymbol{x}}\partial_x + \widehat{\boldsymbol{y}}\partial_y$ & \ref{tide-eddy} \\
$\nabla_3$ & three-dimensional gradient & \ref{tide-eddy} \\
$\nabla$ & gradient & \ref{fn-cstcs}, \ref{act-cstcs}, \ref{lck-drp} \\
\addlinespace
$\widehat{\phantom{u}}\;$ & a Fourier amplitude, as in $\widehat{u}$ & \ref{tide-eddy} \\
$\hat{\phantom{u}}\;$ & a Fourier amplitude, as in $\widehat{\mathbf{u}}_{\mathbf{k}}$ & \ref{trunc-euler} \\
$\overline{\phantom{u}}\;$ & a time average & \ref{tide-eddy} \\
$\langle\cdot\rangle$ & an average & \ref{fn-cstcs}, \ref{act-cstcs}, \ref{lck-drp} \\
$\langle\cdot\rangle_{\rm rms}$ & a root mean square average & \ref{act-cstcs} \\
$\widetilde{\phantom{u}}\;$ & the large-scale part left by the spatial filter & \ref{tide-eddy} \\
$\tilde{\phantom{u}}\;$ & a rescaled or non-dimensional quantity & \ref{act-cstcs} \\
$\dot{(\,)}$ & a derivative with respect to time & \ref{fn-cstcs}, \ref{act-cstcs}, \ref{lck-drp} \\
$\boldsymbol{u}$, bold italic & a vector & all \\
$\mathbf{u}$, bold upright & the Galerkin-truncated field, against the italic untruncated one & \ref{trunc-euler} \\
$\mathbf{S}$, bold upright & a tensor & \ref{act-cstcs} \\
$\mathbb{A}$, blackboard bold & a tensor & \ref{fn-cstcs} \\
$(\,)^*$ & the complex conjugate & \ref{tide-eddy} \\
subscript $g$, $a$ & the geostrophic and the ageostrophic part & \ref{tide-eddy} \\
subscript $s$, $f$ & the slow and the fast part & \ref{tide-eddy} \\
subscript $s$, $b$ & the small and the big droplets & \ref{lck-drp} \\
subscript rms & a root mean square value & \ref{tide-eddy}, \ref{act-cstcs}, \ref{lck-drp} \\

\end{longtable}
\endgroup
   % swap the comment to use the short list

\mainmatter

\subfile{chapters/introduction.tex}
\part{Behaviour of turbulent flows}
\subfile{chapters/internal_tide_eddy.tex}
\subfile{chapters/thermalization.tex}
\part{Behaviour of particles in turbulent flows}
\subfile{chapters/inertial_caustics.tex}
\subfile{chapters/active_caustics.tex}
\subfile{chapters/polydisperse_lucky_drops.tex}
\subfile{chapters/conclusion.tex}
\part*{Appendices}
\phantomsection
\addcontentsline{toc}{part}{Appendices}


\subfile{chapters/appendix.tex}

\backmatter
\printbibliography[heading=bibintoc,title={References}]

\chapter*{Publications included in this thesis}
\addcontentsline{toc}{chapter}{Publications included in this thesis}
\begin{enumerate}[leftmargin=*]
\item \emph{Internal Tide-Eddy Interactions and the Formation of Internal Wave Continuum}. \textbf{Rajarshi, C.}*, Thomas, J.*\\
Under review.

\item \emph{\href{https://doi.org/10.1098/rsta.2025.0024}{The dynamics of thermalisation in the Galerkin-truncated, three-dimensional Euler equation}}. \textbf{Rajarshi, C.}, Khan, M.S., Anand, P., Ray, S.S.*\\
Phil. Trans. R. Soc. A (2026) 384 (2325): 20250024. \href{https://arxiv.org/abs/2512.10788}{ArXiV}.

\item \emph{\href{https://doi.org/10.1103/fwkk-s3w4}{Caustics of finitely dense inertial particles}}. \textbf{Rajarshi, C.}*, Govindarajan, R.\\
Phys. Rev. Fluids 11, 074304 (2026). \href{https://arxiv.org/abs/2601.14179}{ArXiV}.

\item \emph{\href{https://doi.org/10.1103/m3v2-bdvj}{Active caustics}}. Chajwa, R., \textbf{Rajarshi, C.}, Govindarajan, R., Ramaswamy, S.*\\
Phys. Rev. Fluids 11, 033104 (2026) [Editors' Suggestion]. \href{https://doi.org/10.48550/arXiv.2310.01829}{ArXiV}.
\end{enumerate}
\noindent * = corresponding author

\end{document}
